\section{Conclusion}
\label{sec:concl}

% In this work, we introduce Omni-WorldBench, a benchmark for evaluating interaction capabilities in video world models. Unlike existing benchmarks that mainly focus on visual quality or motion realism, Omni-WorldBench emphasizes interaction-driven scene evolution and causal consistency under action prompts. To support this goal, we construct Omni-WorldSuite, a prompt suite covering seven physical principles and task-oriented scenarios, organized by interaction complexity. We further propose Omni-Metric, a unified evaluation framework that measures interaction effect fidelity, non-intervention consistency, spatiotemporal causal coherence, and general visual quality, and aggregates them into an overall AgenticScore. Extensive experiments reveal clear differences in their interaction capabilities. While many models achieve strong visual fidelity and motion smoothness, their ability to maintain causal interaction dynamics remains limited. Our results show that Omni-Metric can effectively capture these differences and aligns well with human judgments. We hope Omni-WorldBench can serve as a standardized testbed for advancing interaction-aware video generation and future world models.

\paragraph{Summary.} In this work, we introduce Omni-WorldBench, the first benchmark dedicated to evaluating the interactive response capabilities of video world models. Unlike existing benchmarks that mainly focus on visual quality or motion realism, Omni-WorldBench emphasizes action-driven scene evolution, intermediate state transitions, and causal consistency under interactive prompts, providing a more comprehensive and holistic evaluation perspective. To support this goal, we establish a rigorous evaluation framework consisting of Omni-WorldSuite, a hierarchical prompt suite spanning diverse interaction levels, physical principles, and task-oriented scenarios, and Omni-Metric, an agent-based evaluation protocol that quantitatively measures the impact of actions on both final outcomes and intermediate state transitions, while also assessing non-intervention consistency, spatiotemporal causal coherence, and visual quality, and aggregating them into an overall AgenticScore. Through a systematic evaluation of 18 video generation models and world models, we reveal substantial gaps between visual realism and true interactivity in current systems: although many models achieve strong visual fidelity and motion smoothness, their ability to maintain causally grounded interaction dynamics remains limited. Our results further show that Omni-Metric can effectively capture these differences. We hope Omni-WorldBench can serve as a standardized testbed for diagnosing current limitations and advancing research on more interactive and causally consistent world models, while being continuously refined and extended through community feedback.


\paragraph{Limitations.}
Despite its broad coverage, Omni-WorldBench still has several limitations. Although Omni-WorldSuite spans diverse physical principles, task-oriented scenarios, and interaction levels, it cannot fully capture the complexity of open-world interactive environments, especially long-horizon and highly dynamic settings. In addition, while Omni-Metric provides a unified protocol for evaluating action-conditioned outcomes and intermediate state transitions, we plan to release human-aligned evaluation results in the future to further complement and validate the assessment of interaction quality. 